% !TeX root = ../main.tex

\chapter{Basic Blocks}

\section{Isolating the Basic Blocks}
% =================================================

Let us forget the coefficients for a moment and isolate the functions

\[
\boxed{
\phi_n(x)=b^{x^n}.
}
\]

Equivalently,

\[
\phi_n(x)=e^{\lambda x^n}.
\]

The first blocks are

\[
\phi_1(x)=b^x,
\]

\[
\phi_2(x)=b^{x^2},
\]

\[
\phi_3(x)=b^{x^3},
\]

\[
\phi_4(x)=b^{x^4},
\]

and so on.

Our series is therefore a linear combination of these blocks:

\[
\boxed{
F(x)=a_0+\sum_{n=1}^{\infty}c_n\phi_n(x).
}
\]

Thus the family

\[
\mathcal{B}
=
\{
1,\phi_1,\phi_2,\phi_3,\ldots
\}
\]

plays a role analogous to the familiar monomial family

\[
\{
1,x,x^2,x^3,\ldots
\}.
\]

However, its algebraic and analytic behaviour is very different.

% =================================================
\section{Connection with Ordinary Monomials}
% =================================================

The blocks are obtained directly from the ordinary monomials.

Let

\[
q_n(x)=x^n.
\]

Then

\[
\phi_n(x)=b^{q_n(x)}.
\]

Therefore

\[
\boxed{
\phi_n=\exp_b\circ q_n,
}
\]

where

\[
\exp_b(y)=b^y.
\]

This observation is important because some of the structure we shall
discover is inherited from the monomials themselves.

For example,

\[
q_n(q_m(x))
=
(x^m)^n
=
x^{mn}.
\]

Thus

\[
\boxed{
q_n\circ q_m=q_{mn}.
}
\]

The monomials already encode multiplication of their indices through
composition.

The exponential blocks preserve this structure after exponentiation.

This distinction will be important:

\begin{quote}
The multiplicative behaviour of the indices is inherited from the
monomials \(x^n\), while the new questions arise from taking linear
combinations of the exponentiated monomials \(b^{x^n}\).
\end{quote}

% =================================================
